\section{Conclusion and Future Outlook}
\label{sec:conclusion}

This paper presented an empirical investigation of Ethereum EIP-4844 blobs as a temporary, verifiable publication layer for structured scientific research artifacts. By designing a deterministic pipeline combining RFC~8785 canonicalization, Zstandard compression, modulus-safe \texttt{FE31} packing, three-tiered hybrid commitments, and a dual-horizon content-addressed archival bridge, we demonstrated that research records can achieve up to a $43.8\times$ cost reduction compared to standard EVM calldata while maintaining $100\%$ bitwise reconstruction parity and sub-second verification latency.

Furthermore, our empirical retrieval study on Ethereum Sepolia highlighted the practical fragility of public consensus Beacon endpoints, establishing that reliable scientific reproducibility strictly requires decoupling consensus publication from durable content-addressed archival persistence. 

As the Ethereum protocol progresses toward PeerDAS (EIP-7594) and increased blob throughput limits (EIP-7691), data availability bandwidth will expand significantly. Future work will investigate decentralized peer-to-peer sidecar dissemination, automated smart contract bounty distribution for archival pinning, and zero-knowledge validity proofs for non-deterministic scientific pipelines.
